\section{Base-model retention}\label{app:base}
We define observed retention as no decrease in any selected benchmark and no increase in held-out NLL. This descriptive rule is not a statistical equivalence test. Table~\ref{tab:nll} reports final NLL; compare recipes within a model, not token-level losses across tokenizers. All final CE and KD variants fail the full retention rule.

\begin{table}[h]
\centering\small
\begin{tabular}{lrrr}
\toprule
Model & Source & CE & CE+KD\\
\midrule
SmolLM2-135M & 2.828018 & 2.856920 & 2.851728\\
SmolLM2-360M & 2.605900 & 2.617682 & 2.614695\\
SmolLM2-1.7B & 2.305988 & 2.326675 & 2.321363\\
Qwen3-0.6B & 2.754853 & 2.781615 & 2.761109\\
Qwen3-1.7B & 2.513415 & 2.531043 & 2.519932\\
\bottomrule
\end{tabular}
\caption{Base-model held-out NLL (nat/token); lower is better.}
\label{tab:nll}
\end{table}

SmolLM2-135M embedding-only adaptation reaches LAMBADA 42.81\% and NLL 3.049, versus source 42.71\% and 2.828. Recovering one metric does not establish overall retention. Two-stage CE reaches LAMBADA 41.90\%; its paired difference from immediate full tuning is +0.80 points [0.29, 1.30], while its difference from the source is $-0.82$ [$-1.59$, $-0.02$]. Intervals are per-metric 95\% intervals over evaluation examples, not training seeds or simultaneous guarantees.

Qwen3-1.7B CE raises ARC-Easy from 68.52\% to 73.57\%, but lowers HellaSwag from 66.37\% to 64.67\%, SciQ from 94.90\% to 93.60\%, and LAMBADA from 63.09\% to 61.69\%. KD partially recovers these losses to 65.59\%, 94.50\%, and 62.25\%, respectively. Base AR continuation and immediate-unfreezing ablations were run for the 135M pilot, not the full model matrix.

\section{Instruction data and ablations}
The instruction model is Qwen3-1.7B, not Qwen3-1.7B-Base. Adaptation uses assistant-only loss on SmolTalk-derived dialogs. Thinking responses are generated by the frozen source; non-thinking responses are corpus targets. Thus CE versus CE+KD isolates additional distribution matching, not the presence versus absence of teacher-generated text. Neither the AR controls nor the mode gap alone establishes a causal explanation for the deficit.

Table~\ref{tab:ablations} reports full IFEval for the SPS arms. Continuation reuses the same corpus; it tests additional updates, not more unique data. Intermediate checkpoints and schedule variants are exploratory comparisons rather than a separately validated checkpoint-selection rule.

\begin{table}[h]
\centering\small
\begin{tabular}{llrr}
\toprule
Recipe & Endpoint & T & NT\\
\midrule
CE & Constant, 512 & 64.14 & 52.50\\
CE+KD & Constant, 512 & 64.33 & 55.82\\
CE & Cosine, 512 & 64.51 & 48.98\\
CE+KD & Cosine, 512 & 62.48 & 59.33\\
CE & Continued, 2048 & 60.81 & 53.79\\
CE+KD & Continued, 2048 & 57.86 & 52.87\\
\bottomrule
\end{tabular}
\caption{SPS instruction-following ablations: full IFEval prompt-level strict accuracy (\%).}
\label{tab:ablations}
\end{table}

\section{Partial MMLU and missing evaluations}
Thinking runs did not all finish. We compare the intersection across the source and twelve adaptations: CE/KD for SPS/AR at constant-LR step 384, constant-LR step 512 and cosine step 512. The 3,859 shared questions cover 25 of 57 subjects, with the last subject incomplete. Question hashes, targets, subjects and generation seeds agree across variants; manifests agree on generation settings and dataset revision. No unfinished JSONL rows were encountered. Truncated generations remain in the denominator and are scored by the extracted final answer; unparseable answers are errors.

\begin{table}[h]
\centering\small
\begin{tabular}{llrr}
\toprule
Schedule/step & Model & Acc. & Invalid\\
\midrule
Source & AR & 67.22 & 19.93\\
Constant/384 & SPS CE & 69.45 & 14.23\\
 & SPS KD & 68.26 & 15.73\\
 & AR CE & 67.27 & 17.34\\
 & AR KD & 68.36 & 16.53\\
Constant/512 & SPS CE & 70.74 & 11.27\\
 & SPS KD & 69.06 & 14.43\\
 & AR CE & 69.94 & 12.80\\
 & AR KD & 67.95 & 16.82\\
Cosine/512 & SPS CE & 67.92 & 16.30\\
 & SPS KD & 69.19 & 16.20\\
 & AR CE & 69.71 & 14.67\\
 & AR KD & 68.54 & 17.05\\
\bottomrule
\end{tabular}
\caption{Thinking MMLU on the identical 3,859-question subset (\%). KD denotes CE+KD.}
\end{table}

There are no saved non-thinking MMLU responses for the source or three constant-step-512 adaptations. The all-variant non-thinking intersection is therefore empty. Nine available variants share 1,377 questions across 11 subjects, but this cannot establish retention against the source. Absolute thinking and non-thinking scores on these different subsets are not comparable. The shared-subset result is descriptive and does not establish statistical significance.

\section{Paired response replacement across generation modes}\label{app:provenance}
We compare continued AR and staged SPS CE under two target conditions. The corpus condition retains the original non-thinking answers; the teacher condition replaces them with frozen-source answers. Thinking examples are identical. Incomplete teacher generations remove the same prompts from both conditions. The controls share 7,376 training prompts, their order, and 461 optimizer steps; we evaluate fixed final checkpoints rather than selecting on benchmark scores. Teacher targets are longer, so supervised-token exposure and the implicit token-normalized mode weighting are not matched. The contrast measures response replacement as a whole, not a length-independent provenance effect.

On all 541 non-thinking IFEval prompts, strict prompt accuracy is 53.79/62.85\% for AR corpus/teacher and 50.65/61.37\% for SPS corpus/teacher, versus 65.43\% for the frozen source. Teacher replacement improves AR by 9.06 percentage points (paired-question bootstrap 95\% interval [4.99,13.12]) and SPS by 10.72 points ([6.84,14.60]). The difference between these effects is 1.66 points ([-2.96,6.28]); a uniquely SPS-specific benefit is not established. SPS with teacher targets remains 4.07 points below the source ([-7.58,-0.74]). These exploratory, per-comparison intervals have no multiplicity correction and do not measure training-seed variability.

Thinking evaluation covers the same 541 prompts. Strict accuracy is 61.74/61.74\% for AR corpus/teacher and 62.66/59.89\% for SPS corpus/teacher, versus 65.62\% for the source. Replacement changes AR by 0.00 points ([-3.88,3.88]) and SPS by -2.77 points ([-6.65,1.11]); these intervals establish neither a thinking loss from replacement nor equivalence. SPS teacher is 5.73 points below the source ([-9.80,-1.66]). The strict architecture-by-replacement interaction is -2.77 points ([-7.95,2.40]); the secondary loose interaction is -5.73 points ([-11.09,-0.37]), an exploratory contrast without multiplicity correction. Malformed-thinking rates rise from 8.87 to 11.28\% for AR and from 7.76 to 12.94\% for SPS. Output-cap rates rise from 14.79 to 16.45\% and from 15.53 to 16.82\%, respectively. These observations motivate a mode-weighting control, but do not identify its causal contribution.

Crossed held-out likelihood exposes target-distribution specialization. Corpus/teacher-trained AR has non-thinking NLL 0.459/1.326 on corpus targets and 0.360/0.135 on teacher targets; SPS has 0.467/1.287 and 0.366/0.146, respectively. Thus better external instruction following coexists with worse corpus-target likelihood. Output-cap rates also decrease from 22.55 to 18.67\% for AR and from 25.14 to 18.67\% for SPS. Broader quality evaluation, repeated seeds and length/token-weighting controls are needed before attributing the gain more narrowly.

\section{Shared-backbone speculative serving pilot}\label{app:serving}
\paragraph{Models and adaptation.}
This experiment is separate from both the base-model conversion matrix and the constant-LR instruction checkpoints in Table~\ref{tab:main}. The AR verifier is the original post-trained \texttt{Qwen/Qwen3-1.7B}; EAGLE-3 \citep{li2025eagle3} uses the same verifier with the published \texttt{AngelSlim/Qwen3-1.7B\_eagle3} draft. For SPS, we initialize from the existing cosine-schedule instruction checkpoint at step 512 and add three prediction embeddings to its first prediction slot. Assistant-masked horizon training first updates only the new embeddings for 102 steps, then the shared backbone and all prediction embeddings for 410 steps. This continuation consumes 8,192 examples, approximately 5.17 million input tokens, including 2.63 million assistant tokens; its context limit is 2,048. Full-stage held-out horizon CE decreases from 2.181 to 1.603. This horizon loss is not a first-slot quality or retention measure.

The final adapted checkpoint supplies both the ordinary SPS baseline and SPS-MTP: there is no separate SPS draft backbone and no inference-time weight change between those two arms. The trained model has four causal prediction slots per real-token prefix; the ordinary serving baseline evaluates only the first slot, whereas SPS-MTP evaluates all four. Later slots propose successors of the pending first-slot sample. Verification processes the candidate real-token block and prediction groups in one backbone traversal. Prediction states can read their own group's earlier prediction slots and causal real-token KV; real states cannot read prediction states, and only real-token KV persists. SPS proposals are greedy; target sampling uses vLLM's rejection sampler. The AR and SPS verifier weights differ, so their relative speedups must be interpreted against separate baselines.

\paragraph{Evaluation and timing.}
We run native vLLM \citep{kwon2023vllm} on one A100-SXM4-80GB per evaluation job, BF16, tensor parallelism one, with a single active request and Triton attention. A seeded, subject-balanced round-robin procedure selects the same 100 STEM test questions for every arm and mode. This is a small fixed pilot, not full MMLU, a uniformly sampled MMLU test, or the 3,859-question intersection in the instruction analysis. Thinking uses temperature 0.6, top-$p$ 0.95, top-$k$ 20 and a 2,048-token output cap; non-thinking uses 0.7, 0.8, 20 and 512, respectively. All speculative arms propose up to three tokens. Prefix caching, chunked prefill and asynchronous scheduling are disabled. Two requests with shortened output caps warm the engine before three timed repetitions of the same questions and per-question seeds. Accuracy is reported from the first timed repetition; timing repeats are not independent quality samples.

For repetition $r$, aggregate decode throughput is
\[
D_r=\frac{\sum_{i=1}^{100}\max(n_{ir}-1,0)}{\sum_{i=1}^{100}(t^{\rm last}_{ir}-t^{\rm first}_{ir})},
\]
where $n_{ir}$ counts generated tokens and vLLM timestamps delimit the first and last generated token events. Prefill is excluded. Tables report $\mathrm{median}_r D_r$ and the median paired ratio $D^{\rm spec}_r/D^{\rm verifier}_r$. The reported speedup need not equal the ratio of the two separately reported throughput medians. Request latency is measured separately around generation with device synchronization. Neither throughput ratios nor request latency compare identical output lengths across draft methods: sampling can produce different continuations. Invalid final answers count as errors, and capped outputs remain in the accuracy denominator.

\paragraph{CUDA graph scope and correctness.}
We compare eager execution with full \emph{target decode} CUDA graphs, without Torch compilation; prefill is outside the graph. The graph implementation also replaces SPS per-forward Python state assignments with copies into persistent hidden and position buffers, which graph replay updates. Eager and graph runs thus use separately pinned serving implementations with identical model weights. Named worker-extension RPC methods count actual target replay during initial untimed requests, and restore the replay method before timed generation. Every graph arm records positive target replay evidence; this does not establish capture of the separate EAGLE drafter or imply that every irregular decode block uses a graph. Both SPS arms pass exact greedy parity with the original training-runtime reference on two prefixes of up to 16 generated tokens. Across the timed suite, all 2,400 graph-generated token sequences exactly match their corresponding eager sequences. These checks do not establish general sampling-distribution equivalence between speculative and non-speculative generation.

\paragraph{Expanded correctness diagnostics.}
A subsequent audit passed known-distribution rejection-kernel checks but found greedy SPS/SPS-MTP differences in 30 of 112 cases, while eager and graph execution agreed within each arm. Cases reuse eight prefixes across output limits and EOS settings; they are not independent quality samples. Matching both arms to four slots reduced differences to 16/112 with BF16, then 6/112 with an FP32 output head. An FP16-backbone control retaining the FP32 head left 2/112 differences, both EOS settings of one prefix distinct from the three affected prefixes in the preceding control. This last intervention changes both weight rounding and arithmetic from the same FP32 export. Matched four-slot full-prefix distributions were identical between ordinary and speculative execution on all 13 diagnostic histories in each precision condition, but this did not imply cached equivalence.

Tracing the three first-divergence histories from the BF16-backbone/FP32-head control reproduced the earlier outputs. Every observed greedy selection followed its verifier argmax; the sequential/speculative target distributions already differed before rejection sampling, with total variation 0.0083--0.0152 on those histories. Cold full-prefix argmax agreed with speculative decoding on all three. These results locate the observed discrepancies upstream of rejection sampling and demonstrate precision sensitivity; they do not prove cache correctness or general positive-temperature sampling equivalence. The original correctness gate remains failed. These diagnostic precision changes were not used for the throughput or quality tables below.

\begin{table*}[t]
\centering\small
\begin{tabular}{lrrrrrr}
\toprule
& \multicolumn{3}{c}{Thinking} & \multicolumn{3}{c}{Non-thinking}\\
Model & Acc. (\%) & Decode tok/s & Speedup & Acc. (\%) & Decode tok/s & Speedup\\
\midrule
AR & 62 & 193.55 & 1.00$\times$ & 57 & 199.48 & 1.00$\times$\\
AR + EAGLE-3 & 62 & 282.44 & 1.46$\times$ & 63 & 319.62 & 1.60$\times$\\
SPS & 70 & 161.54 & 1.00$\times$ & 55 & 185.60 & 1.00$\times$\\
SPS + MTP & 63 & 254.71 & 1.58$\times$ & 59 & 334.87 & 1.80$\times$\\
\bottomrule
\end{tabular}
\caption{Fixed MMLU STEM pilot with full target decode CUDA graphs. Speedup uses the matching verifier with graphs: AR for EAGLE-3, final SPS for SPS-MTP. Rows therefore do not share one verifier. Accuracy is based on 100 questions; three timing repetitions reuse seeds.}
\label{tab:serving-graphs}
\end{table*}

\begin{table*}[t]
\centering\small
\begin{tabular}{lrrrrrr}
\toprule
& \multicolumn{3}{c}{Thinking} & \multicolumn{3}{c}{Non-thinking}\\
Model & Eager tok/s & Eager speedup & Graph speedup & Eager tok/s & Eager speedup & Graph speedup\\
\midrule
AR & 43.03 & 1.00$\times$ & 1.00$\times$ & 44.99 & 1.00$\times$ & 1.00$\times$\\
AR + EAGLE-3 & 68.38 & 1.59$\times$ & 1.46$\times$ & 79.51 & 1.75$\times$ & 1.60$\times$\\
SPS & 36.66 & 1.00$\times$ & 1.00$\times$ & 37.19 & 1.00$\times$ & 1.00$\times$\\
SPS + MTP & 74.67 & 2.04$\times$ & 1.58$\times$ & 85.95 & 2.32$\times$ & 1.80$\times$\\
\bottomrule
\end{tabular}
\caption{Execution-mode sensitivity on the same pilot. Eager and graph speedups each use their own execution-mode-matched verifier baseline. Graph throughput is in Table~\ref{tab:serving-graphs}; answers and accuracy are unchanged between execution modes.}
\label{tab:serving-eager}
\end{table*}

\paragraph{Results and interpretation.}
Tables~\ref{tab:serving-graphs}--\ref{tab:serving-eager} show that graphs increase absolute throughput for every arm but reduce relative speculative gains. With graphs, SPS-MTP achieves 1.58$\times$ thinking and 1.80$\times$ non-thinking speedups; observed timing-repeat ranges are 1.57--1.58$\times$ and 1.80--1.81$\times$. These are timing ranges, not confidence intervals over questions or training seeds. EAGLE-3 reaches 1.46$\times$ and 1.60$\times$. Its absolute thinking throughput exceeds SPS-MTP (282.44 versus 254.71 tok/s), while SPS-MTP has higher absolute non-thinking throughput (334.87 versus 319.62 tok/s). For SPS/SPS-MTP, the median over repetitions of mean request generation latency is 6.259/4.106 seconds in thinking and 0.301/0.168 seconds in non-thinking; request latency also depends on the sampled response lengths.

Faster non-thinking decoding co-occurs with shorter outputs and higher draft acceptance. SPS-MTP averages 1,038 thinking versus 47 non-thinking output tokens, with 38.9\% versus 52.7\% of proposed tokens accepted. EAGLE-3 acceptance is 33.5\% versus 41.5\%. Acceptance counters include warmup and, for SPS arms, the untimed greedy checks; we did not separately profile the causal contributions of acceptance, context length and launch overhead.

No quality-preserving acceleration claim follows from this pilot. SPS-MTP thinking accuracy is 63\%, versus 70\% for the same SPS verifier without speculation; its non-thinking accuracy is 59\%, versus 55\%. For AR/EAGLE-3 the corresponding accuracies are 62/62\% and 57/63\%. Thinking output-cap rates are 13/18\% for SPS/SPS-MTP and 24/26\% for AR/EAGLE-3; non-thinking cap rates are 1/1\% and 7/5\%, respectively. Mode-dependent sampling, a small fixed question set and output truncation limit the quality interpretation. A fixed-length decode study, repeated quality seeds, broader tasks and a same-verifier comparison of draft methods remain separate follow-up experiments.

\section{Fixed-weight DFlash mask ablation}\label{app:dflash-mask}
We use the published five-layer DFlash-b16 drafter \citep{chen2026dflash} with
the original \texttt{Qwen/Qwen3-4B} as the verifier. Both masks share
weights, conditioning layers, the verifier embedding and output head,
and the procedure for caching draft K/V per layer. For each block of $k+1$ positions, the verifier supplies
the first token and the drafter predicts $k$ proposals in one forward pass.
The original DFlash mask allows attention to the entire block. The
SPS-style mask lets each query see the verified prefix and the
block positions through its own. We test $k\in\{4,8\}$.

We use 24 fixed MMLU STEM prompts with thinking enabled, greedy decoding
and a cap of 128 new tokens. Two timing repeats per prompt rotate the order
of the four configurations. All run in BF16 with PyTorch SDPA and one
request at a time on one A100-SXM4-80GB. Decode throughput pools the tokens
and wall time after prefill and the first token. Acceptance is the fraction
of proposals accepted in the contiguous prefix. We compute acceptance and
$\tau$ (mean accepted proposals plus one) over fully emitted verification
rounds. We compare each output with sequential greedy decoding.

\begin{table}[t]
\centering
\small
\setlength{\tabcolsep}{4pt}
\begin{tabular}{lrrrrr}
\toprule
Mask & $k$ & Tok/s & $\tau$ & Accept. & First \\
\midrule
Bidirectional & 4 & 45.78 & 2.31 & 32.8 & 60.7 \\
Causal & 4 & 42.03 & 2.12 & 28.1 & 54.6 \\
\addlinespace
Bidirectional & 8 & 49.06 & 2.46 & 18.3 & 60.7 \\
Causal & 8 & 43.87 & 2.21 & 15.1 & 53.6 \\
\bottomrule
\end{tabular}
\caption{Fixed-weight DFlash mask ablation. Accept. is the percentage of
proposals accepted. First is the percentage of rounds accepting the first
proposal. Both statistics use fully emitted verification rounds.}
\label{tab:dflash-mask}
\end{table}

Table~\ref{tab:dflash-mask} shows lower acceptance with the causal mask
at both widths, including on the first proposal. Decode throughput falls
by 8.2\% at $k=4$ and 10.6\% at $k=8$. GPU timings per round
remain close across masks, at roughly 7.2 ms for drafting and 43 ms for
verification. The smaller $\tau$ leaves more rounds to complete the same
128-token outputs. Moving from four to eight proposals raises $\tau$ and
throughput under both masks while lowering acceptance as a fraction of
proposals.

For each width, the two masks emit identical token sequences on all 24
prompts, also across repeats. Compared with sequential greedy Qwen3-4B,
11 of 24 outputs match at $k=4$ and 12 of 24 at $k=8$ for both masks.
The comparison therefore measures acceptance and timing on matched output
sequences within each width. DFlash-b16 was trained with bidirectional
blocks of size 16. Whether training with the causal mask closes the gap
remains open.

\section{Evidence and draft scope}
The code registries are the single source of experimental configuration. Base runs use \texttt{scripts/ar\_sps\_config.py}; instruction, scheduling and continuation runs use the corresponding \texttt{instruct}, \texttt{sched} and \texttt{extend} configuration modules. Serving definitions, training, timing and aggregation use the \texttt{vllm\_sps} and \texttt{serving\_graphs} script families. Training and evaluation workers are separate. The accompanying draft README and Russian appendix notes list source paths, measured evidence locations and pinned serving revisions.

This draft includes the shared-backbone vLLM MTP pilot, but excludes the earlier alternative speculative prototypes and the long base-training follow-up. The remaining full benchmark tables, complete model/data bibliography and publication-ready uncertainty analysis are not yet incorporated. The expanded correctness diagnostics are follow-up evaluation jobs; the serving throughput and quality tables retain the original BF16 measurements.
